\documentclass[11pt,a4paper]{article}

\usepackage[utf8]{inputenc}
\usepackage[T1]{fontenc}
\usepackage[french]{babel}
\usepackage{lmodern}
\usepackage{microtype}
\usepackage{amsmath,amssymb,amsfonts}
\usepackage{geometry}
\usepackage{booktabs}
\usepackage{hyperref}
\usepackage{xcolor}
\usepackage{enumitem}

\geometry{margin=2.4cm}
\hypersetup{
  colorlinks=true,
  linkcolor=blue!60!black,
  citecolor=blue!60!black,
  urlcolor=blue!60!black,
  pdftitle={Conclusions cosmologiques d'un modele GUP UV/IR},
  pdfauthor={Notes de travail}
}
\setlength{\parindent}{0pt}
\setlength{\parskip}{0.6em}
\setlist[itemize]{topsep=3pt,itemsep=2pt,parsep=2pt}
\setlist[enumerate]{topsep=3pt,itemsep=2pt,parsep=2pt}

\newcommand{\rhoP}{\rho_{\rm P}}
\newcommand{\ellP}{\ell_{\rm P}}
\newcommand{\rhoDE}{\rho_{\rm DE}}
\newcommand{\rhoReg}{\rho_{\rm reg}}
\newcommand{\rhoBH}{\rho_{\rm BH}}
\newcommand{\LH}{L_H}
\newcommand{\EP}{E_{\rm P}}
\newcommand{\EH}{E_H}
\newcommand{\EDE}{E_{\rm DE}}
\newcommand{\mpl}{m_{\rm P}}

\title{Conclusions cosmologiques d'un cadre GUP--UV/IR\\
\large Energie noire, borne d'horizon et projection holographique}
\author{Notes de travail}
\date{22 mai 2026}

\begin{document}
\maketitle

\begin{abstract}
On examine les conclusions cosmologiques d'un modele ou une cellule de phase de
type GUP regularise l'ultraviolet, tandis qu'une contrainte gravitationnelle
globale impose une borne infrarouge de type horizon. Le resultat central est
negatif et constructif : la regularisation GUP seule transforme la divergence
de l'energie du vide en une valeur finie mais encore planckienne. L'echelle
observee de l'energie noire apparait seulement lorsqu'une projection
gravitationnelle ou holographique relie le cutoff UV a une longueur IR
cosmologique. Le probleme cosmologique est donc essentiellement UV/IR.
\end{abstract}

\section{Constantes et echelles}

Les echelles de Planck sont
\begin{align}
\ellP &= \sqrt{\frac{\hbar G}{c^3}},&
\EP &= \sqrt{\frac{\hbar c^5}{G}},&
\rhoP &= \frac{c^7}{\hbar G^2}.
\end{align}
L'echelle de Hubble est
\begin{equation}
\LH=\frac{c}{H_0}.
\end{equation}
La densite critique actuelle, exprimee comme densite d'energie, vaut
\begin{equation}
\rho_c=\frac{3c^2H_0^2}{8\pi G},
\end{equation}
et
\begin{equation}
\rhoDE=\Omega_\Lambda\rho_c.
\end{equation}

\section{Resultat negatif : le GUP seul ne suffit pas}

La mesure de phase effective est
\begin{equation}
dN=\frac{d^3x\,d^3p}{h^3(1+\beta p^2)^3},
\qquad
\beta=\beta_0\frac{\ellP^2}{\hbar^2}.
\end{equation}
L'integrale de point zero massless donne
\begin{equation}
\boxed{
\rhoReg=\frac{g\rhoP}{16\pi^2\beta_0^2}.
}
\label{eq:rho_reg}
\end{equation}
Si $\beta_0\sim1$ et $g\sim1$, alors $\rhoReg$ est finie mais d'ordre
planckien, donc beaucoup trop grande devant l'energie noire observee :
\begin{equation}
\rhoDE\sim10^{-123}\rhoP.
\end{equation}

\begin{equation}
\boxed{
\text{GUP seul : }\infty\longrightarrow \text{valeur finie mais trop grande.}
}
\end{equation}

\section{Projection horizon--Planck}

La relation cosmologique exacte est
\begin{equation}
\boxed{
\rhoDE
=
\Omega_\Lambda\rho_c
=
\frac{3\Omega_\Lambda}{8\pi}\rhoP
\left(\frac{\ellP}{\LH}\right)^2.
}
\label{eq:rho_de_uvir}
\end{equation}
Le facteur minuscule vient donc du rapport d'aires
\begin{equation}
\left(\frac{\ellP}{\LH}\right)^2.
\end{equation}
L'energie noire ressemble a une densite planckienne projetee par une surface
cosmologique d'horizon.

\section{Borne trou noir}

Une sphere de rayon $L$ contenant une densite d'energie $\rho$ a
\begin{equation}
E=\rho\frac{4\pi}{3}L^3.
\end{equation}
Un trou noir de rayon de Schwarzschild $L$ a
\begin{equation}
E_{\rm BH}=\frac{c^4L}{2G}.
\end{equation}
La condition $E\leq E_{\rm BH}$ donne
\begin{equation}
\boxed{
\rho\leq\rhoBH(L)=\frac{3c^4}{8\pi GL^2}
=
\frac{3}{8\pi}\rhoP\left(\frac{\ellP}{L}\right)^2.
}
\label{eq:rho_bh}
\end{equation}
Pour $L=\LH$, on retrouve
\begin{equation}
\rhoBH(\LH)=\rho_c,
\qquad
\rhoDE=\Omega_\Lambda\rhoBH(\LH).
\end{equation}

\section{Densite effective UV/IR}

La formule cosmologiquement pertinente n'est pas simplement $\rhoReg$. Une
formulation minimale est
\begin{equation}
\boxed{
\rho_{\rm eff}(L)
=
\min\left[
\rhoReg,
\rhoBH(L)
\right].
}
\label{eq:rho_eff_min}
\end{equation}
Donc
\begin{equation}
\boxed{
\rho_{\rm eff}(L)
=
\min\left[
\frac{g\rhoP}{16\pi^2\beta_0^2},
\frac{3c^4}{8\pi GL^2}
\right].
}
\end{equation}
Dans l'univers actuel, $\LH\gg\ellP$, donc la branche active est la branche
gravitationnelle :
\begin{equation}
\rho_{\rm eff}(\LH)=\rhoBH(\LH)=\rho_c.
\end{equation}
L'energie noire observee correspond ensuite a une fraction
\begin{equation}
\rhoDE=\Omega_\Lambda\rho_c.
\end{equation}

\section{Facteur de projection}

Sur la branche gravitationnelle,
\begin{equation}
\rho_{\rm grav}(L)=\mathcal{F}(L)\rhoReg.
\end{equation}
Ainsi
\begin{equation}
\mathcal{F}(L)=\frac{\rhoBH(L)}{\rhoReg}
=
\boxed{
\frac{6\pi\beta_0^2}{g}
\left(\frac{\ellP}{L}\right)^2
}.
\label{eq:projection_factor}
\end{equation}
Le facteur $10^{-122}$ n'est donc pas introduit arbitrairement : il vient de la
comparaison entre l'echelle de Planck et l'echelle d'horizon.

\section{Interpretation holographique}

Un volume de taille $L$ contient naivement
\begin{equation}
N_{\rm vol}\sim\left(\frac{L}{\ellP}\right)^3
\end{equation}
cellules planckiennes. La borne de Bekenstein--Hawking donne plutot
\begin{equation}
\frac{S_{\rm BH}}{k_B}
=
\frac{A}{4\ellP^2}
=
\pi\left(\frac{L}{\ellP}\right)^2.
\end{equation}
Donc
\begin{equation}
\boxed{
\frac{S_{\rm BH}}{k_B}\sim N_{\rm vol}^{2/3}.
}
\end{equation}
Le GUP donne une regularisation volumique ; la gravite projette en loi d'aire.

\section{Echelle d'energie noire}

Definissons
\begin{equation}
\EDE=\left[\rhoDE(\hbar c)^3\right]^{1/4}.
\end{equation}
Avec
\begin{equation}
\EP=\sqrt{\frac{\hbar c^5}{G}},
\qquad
\EH=\hbar H_0,
\end{equation}
on obtient
\begin{equation}
\boxed{
\EDE=
\left(\frac{3\Omega_\Lambda}{8\pi}\right)^{1/4}
\sqrt{\EP\EH}.
}
\label{eq:e_de}
\end{equation}
L'echelle meV de l'energie noire est la moyenne geometrique entre l'echelle de
Planck et l'echelle de Hubble.

\section{Equation d'etat}

Obtenir la bonne densite ne suffit pas. Une vraie energie noire doit avoir
\begin{equation}
w=\frac{p}{\rho}\simeq-1.
\end{equation}
Une somme naive de modes massless regularises peut ressembler a du rayonnement,
$p=\rho/3$, tandis qu'une constante cosmologique correspond a
\begin{equation}
T_{\mu\nu}^{\rm vac}=-\rho_{\rm vac}g_{\mu\nu}.
\end{equation}
La formulation correcte doit donc etre covariante et definir la partie
gravitationnellement active du vide, pas seulement l'integrale de modes.

\section{Choix de la longueur IR}

Une densite du type
\begin{equation}
\rho_X(L)=\eta\frac{3c^4}{8\pi GL^2}
\end{equation}
depend du choix de $L$. Si $L=c/H$, alors $\rho_X\propto H^2$ et le terme peut
se comporter comme une renormalisation de $G$.

Dans les modeles d'energie noire holographique, on prend souvent l'horizon des
evenements futur
\begin{equation}
L=R_{\rm eh}=a(t)\int_t^\infty\frac{c\,dt'}{a(t')}.
\end{equation}
Dans le modele plat non interactif standard, l'equation d'etat devient
\begin{equation}
\boxed{
w_{\rm HDE}
=
-\frac13-\frac{2}{3c_{\rm hde}}\sqrt{\Omega_{\rm DE}},
}
\label{eq:w_hde}
\end{equation}
ou $c_{\rm hde}$ est le parametre sans dimension du modele holographique. Cette
relation montre que la dynamique de $L$ est le coeur physique du modele.

\section{Transition primordiale}

Dans l'univers primordial, la borne $\rhoBH(L)$ augmente lorsque $L$ diminue.
La transition entre branche GUP et branche horizon est definie par
\begin{equation}
\rhoReg=\rhoBH(L_*).
\end{equation}
Donc
\begin{equation}
\frac{g\rhoP}{16\pi^2\beta_0^2}
=
\frac{3}{8\pi}\rhoP\left(\frac{\ellP}{L_*}\right)^2,
\end{equation}
et
\begin{equation}
\boxed{
L_*=\beta_0\sqrt{\frac{6\pi}{g}}\,\ellP.
}
\label{eq:l_star}
\end{equation}
Pour $\beta_0\sim g\sim1$, cette transition est proche de l'echelle de
Planck. Le GUP est alors surtout une physique du tres jeune univers, tandis que
l'energie noire actuelle est controlee par l'horizon cosmologique.

\section{Predictions et tests naturels}

Un modele cosmologique issu de cette architecture devrait produire au moins une
des signatures suivantes :
\begin{enumerate}
 \item une energie noire legerement dynamique, avec $w(a)$ proche de $-1$ ;
 \item une relation entre $\rhoDE$ et une longueur d'horizon $L(t)$ ;
 \item une transition UV/IR dans l'univers primordial ;
 \item l'absence d'une longueur minimale effective macroscopique, car
 $\beta_0$ n'a plus besoin d'etre ajuste a $10^{60}$ ;
 \item une formulation covariante produisant $T_{\mu\nu}\simeq-\rho g_{\mu\nu}$.
\end{enumerate}

\section{Programme theorique}

Pour passer d'une architecture plausible a une theorie cosmologique, il faut :
\begin{enumerate}
 \item deriver rigoureusement la mesure GUP depuis l'algebre deformee ;
 \item calculer la densite de vide regularisee avec les bons signes
 bosons--fermions ;
 \item soustraire ou neutraliser le vide plat de facon covariante ;
 \item montrer que la partie gravitationnellement active est bornee par
 l'horizon ;
 \item determiner dynamiquement le bon choix de $L$ ;
 \item obtenir une equation d'etat proche de $-1$ ;
 \item verifier la compatibilite avec les donnees cosmologiques.
\end{enumerate}

\section{Conclusion}

La conclusion centrale est
\begin{equation}
\boxed{
\text{l'energie noire n'est pas la densite brute du vide GUP.}
}
\end{equation}
Elle serait plutot
\begin{equation}
\boxed{
\text{le residu gravitationnel ou holographique d'un vide UV regularise.}
}
\end{equation}
La chaine logique devient
\begin{equation}
\boxed{
\text{GUP}\Rightarrow\rhoReg<\infty,
\qquad
\text{gravite+horizon}\Rightarrow
\rho_{\rm grav}\sim\rhoP\left(\frac{\ellP}{\LH}\right)^2.
}
\end{equation}
Autrement dit, la constante cosmologique est probablement un probleme UV/IR,
pas seulement UV.

\begin{thebibliography}{9}
\bibitem{Planck2018}
N. Aghanim et al. (Planck Collaboration),
``Planck 2018 results. VI. Cosmological parameters'',
\emph{Astronomy \& Astrophysics} \textbf{641}, A6 (2020),
arXiv:1807.06209.

\bibitem{NISTCODATA}
E. Tiesinga, P. J. Mohr, D. B. Newell and B. N. Taylor,
``CODATA recommended values of the fundamental physical constants: 2022'',
\emph{Reviews of Modern Physics} \textbf{97}, 025002 (2025).

\bibitem{CKN}
A. Cohen, D. Kaplan and A. Nelson,
``Effective Field Theory, Black Holes, and the Cosmological Constant'',
\emph{Physical Review Letters} \textbf{82}, 4971 (1999),
arXiv:hep-th/9803132.

\bibitem{LiHDE}
M. Li,
``A Model of Holographic Dark Energy'',
\emph{Physics Letters B} \textbf{603}, 1--5 (2004),
arXiv:hep-th/0403127.

\bibitem{Kempf}
A. Kempf, G. Mangano and R. B. Mann,
``Hilbert space representation of the minimal length uncertainty relation'',
\emph{Physical Review D} \textbf{52}, 1108 (1995),
arXiv:hep-th/9412167.

\bibitem{Hossenfelder}
S. Hossenfelder,
``Minimal Length Scale Scenarios for Quantum Gravity'',
\emph{Living Reviews in Relativity} \textbf{16}, 2 (2013).
\end{thebibliography}

\end{document}
